\section{系统设计与实现}
\label{sec:design}

RucksFS 是一个用 Rust 实现的分布式文件系统原型。系统由三部分组成：FUSE 客户端提供 POSIX 入口，MetadataServer 负责命名空间和 inode 元数据，DataServer 负责文件数据的读写，客户端和服务端之间走 gRPC。整个设计围着元数据访问路径转，重点讨论三件事：目录项怎么映射到 RocksDB 键空间，目录修改在并发下如何保持正确性，以及客户端怎么在元数据服务和数据服务之间做协调。

\subsection{设计目标}

在动手实现之前，先明确几个从一开始就希望系统具备的性质，它们会在后续章节反复影响具体决策。

最基本的一条是\textbf{标准 POSIX 访问入口}。应用程序应该能用普通的 \texttt{open()}、\texttt{rename()}、\texttt{unlink()} 等系统调用访问挂载点，不需要为 RucksFS 改写任何接口。这也是选择 FUSE 作为对外接口层的直接原因。

接下来是\textbf{元数据与文件数据在服务层的显式分离}。一方面，二者的访问模式本来就不一样：元数据偏随机、偏小写入，和事务/并发控制耦合度高；文件数据则更接近顺序 I/O，路径更直。另一方面，本文真正想研究的就是元数据路径，把数据路径单独拎出去，也能避免数据块管理的细节干扰对元数据组织方式的讨论。

第三是\textbf{并发下的正确性要经得起检查}。名字唯一、目录非空判断、跨目录重命名的原子性、延迟删除这些语义，不能寄希望于“多数情况下不会冲突”这种经验假设，需要用事务边界把它们明确固定下来。这一点是选择 RocksDB 悲观事务的最直接动机。

最后是\textbf{把性能瓶颈集中到少数可分析的关键路径上}，例如目录项查找、事务提交与父目录属性更新。这样即便某个地方成了瓶颈，也好定位和针对性优化，而不至于把文件系统语义整体外包给通用数据库之后，再在黑盒后端里追查性能问题。

\subsection{系统整体架构}
\label{sec:architecture}

RucksFS 的运行时由三个角色组成：挂在客户端机器上的 FUSE 客户端、负责命名空间和 inode 状态的 MetadataServer，以及负责文件内容读写的 DataServer。客户端收到内核的 FUSE 请求后，不会一股脑都转给同一个后端，而是先过一下本地的 \texttt{VfsCore} 路由层，再按请求类型决定应该找元数据面还是数据面。图~\ref{fig:overall_arch} 画了整体的请求路径。

\begin{figure}[htbp]
    \centering
    \begin{tikzpicture}[
        font=\small,
        box/.style={draw, rounded corners=2pt, fill=black!5, align=center, minimum width=2.45cm, minimum height=0.95cm},
        arr/.style={->, >=stealth, thick}
    ]
    \node[box] (client) at (0,0) {FUSE Client\\\texttt{VfsCore}};
    \node[box] (mds) at (3.4,1.35) {MetadataServer\\元数据操作\\\scriptsize \texttt{lookup/create/rename}};
    \node[box] (ds) at (3.4,-1.35) {DataServer\\数据操作\\\scriptsize \texttt{read/write/truncate}};
    \node[box] (rocks) at (6.7,1.35) {RocksDB\\元数据};
    \node[box] (raw) at (6.7,-1.35) {RawDiskDataStore\\\texttt{data.raw}};

    \draw[arr] (client.east) -- (mds.west);
    \draw[arr] (client.east) -- (ds.west);
    \draw[arr] (mds.east) -- (rocks.west);
    \draw[arr] (ds.east) -- (raw.west);
    \end{tikzpicture}
    \caption{RucksFS 的整体架构与主请求路径}
    \label{fig:overall_arch}
\end{figure}

从请求走向上看，这种拆分对应三类路径。

第一类是元数据路径。\texttt{lookup}、\texttt{create}、\texttt{mkdir}、\texttt{unlink}、\texttt{rename}、\texttt{readdir} 这些请求经 FUSE 客户端进 \texttt{VfsCore}，然后被转给 MetadataServer，再由它去访问 RocksDB。

第二类是数据路径。\texttt{read}、\texttt{write}、\texttt{truncate}、\texttt{flush} 这些请求由 \texttt{VfsCore} 选对应的 DataServer，交给原始磁盘后端执行。

第三类是协调路径。有些操作天然跨元数据面和数据面。比如客户端写完数据后要通知 MetadataServer 更新文件大小和 \texttt{mtime}；又比如 \texttt{unlink} 或覆盖式 \texttt{rename} 完成后，MetadataServer 会在响应结构里告诉客户端哪些 inode 需要在 DataServer 侧清理数据。

客户端不直接持有 RocksDB，也不直接管原始磁盘布局。它只在 FUSE 请求、MetadataServer、DataServer 之间做语义编排。这样做的直接好处是，文件系统的正确性主要集中在元数据事务设计和客户端协调逻辑里，边界比较清楚，调试路径也更可控。

整体布局定下来之后，后面几节就按“元数据在 RocksDB 键空间里怎么组织”、“事务模型如何保证并发正确性”、“核心 POSIX 操作怎么映射到键空间和事务上”、“客户端如何协调元数据面和数据面”这个顺序展开。

\subsection{元数据模型与存储布局}
\label{sec:datamodel}

这一节先给出面向实现的元数据全景，再分别展开列族划分、键编码和值编码上的细节。“全景”这一段只负责回答三个问题：系统保存哪些对象，这些对象分别落在哪个列族，以及它们的 key/value 编码是什么样子。至于为什么这么划分、这种编码对各类操作意味着什么，留到后面小节再展开。

\subsubsection{元数据方案总览}

RucksFS 从一棵目录树出发理解命名空间。用户看到的是按目录组织的层次结构，但真正需要持久化的核心对象只有两类：inode 和 dentry（目录项）。inode 代表文件系统对象本体，存类型、权限、大小、链接计数、时间戳这些属性；dentry 代表目录树上一条父子边，记“某个父目录下的某个名字指向哪个 inode”。普通文件、目录、符号链接都先由 inode 标识，再通过 dentry 挂到目录树里。

这样一来，路径本身不是单独保存的对象，而是靠一串 dentry 逐级串出来的。以 \texttt{/docs/paper.tex} 为例，先在根目录下找 \texttt{docs} 对应的 dentry，拿到目录 inode；再在那个目录下找 \texttt{paper.tex} 对应的 dentry，最终定位到目标 inode。围绕这两个核心对象，RucksFS 在 RocksDB 里分别组织 inode 记录和 dentry 记录，另外再补上数据位置、属性增量和系统状态等辅助元数据。图~\ref{fig:cf_kv_layout} 给出了这套方案的整体视图。

\begin{figure}[htbp]
    \centering
    \begin{tikzpicture}[
        font=\small,
        cf/.style={draw, rounded corners=2pt, fill=black!3, align=left, text width=11.8cm, inner sep=6pt},
        title/.style={font=\small\bfseries}
    ]
    \node[cf] (cf1) at (0,0) {\textbf{\texttt{inodes} 列族}\\[2pt]
        inode 属性记录：\texttt{[I][inode: u64 BE]} $\rightarrow$ 57 字节定长 \texttt{InodeValue}\\
        数据位置记录：\texttt{[L][inode: u64 BE]} $\rightarrow$ \texttt{DataLocation}\\
        符号链接记录：\texttt{[S][inode: u64 BE]} $\rightarrow$ 目标路径字符串};

    \node[cf, below=6mm of cf1] (cf2) {\textbf{\texttt{dir\_entries} 列族}\\[2pt]
        目录项记录：\texttt{[D][parent: u64 BE][name: UTF-8]}\\
        \hspace*{4.5em}$\rightarrow$ \texttt{[child\_inode: u64 BE][child\_mode: u32 BE]}};

    \node[cf, below=6mm of cf2] (cf3) {\textbf{\texttt{delta\_entries} 列族}\\[2pt]
        属性增量记录：\texttt{[X][inode: u64 BE][seq: u64 BE]} $\rightarrow$ \texttt{DeltaOp}};

    \node[cf, below=6mm of cf3] (cf4) {\textbf{\texttt{system} 列族}\\[2pt]
        系统状态记录：系统保留键 $\rightarrow$ 分配器状态等元信息};
    \end{tikzpicture}
    \caption{RucksFS 元数据记录与列族布局总览}
    \label{fig:cf_kv_layout}
\end{figure}

图~\ref{fig:cf_kv_layout} 可以按列族来读：\texttt{inodes} 中集中保存 inode 相关记录，\texttt{dir\_entries} 单独保存目录项，\texttt{delta\_entries} 保存属性增量，\texttt{system} 保存低频系统状态。这样先把“有哪些记录、分别放在哪里”看清楚，后续再展开各类记录的编码细节与设计考量。

\subsubsection{列族划分与设计考量}

有了上面的全景布局，下一个问题就是：为什么这些对象要被这么分组，而不是全扔进一个统一键空间？本文的做法是按访问模式把 RocksDB 实例拆成 4 个列族——\texttt{inodes}、\texttt{dir\_entries}、\texttt{delta\_entries}、\texttt{system}——各自承担 inode 点查、目录项存储、属性增量追加、系统状态保存这几类不同的任务。

\texttt{inodes} 列族承担“对象本体”的点查负载。inode 基础属性、\texttt{DataLocation} 和符号链接目标语义上并不完全同类，但访问模式一致——都是按 inode 直接定位的点查——所以放在同一列族里。这样就不用为少量辅助元数据再单独维护一套 MemTable 和压缩状态。

\texttt{dir\_entries} 列族承担命名空间的边关系。\texttt{lookup} 要按“父 inode + 名字”定位单条目录项，\texttt{readdir} 要按父目录枚举全部子项，这两种访问都依赖前缀局部性。所以这个列族启用了 9 字节的前缀提取器，让同一父目录下的目录项按前缀连续排列。相比“目录项和 inode 混在一起”，它更利于目录遍历和前缀 Bloom 过滤器生效；相比“用完整路径做键”，它又避免了跨目录 \texttt{rename} 要对整棵子树改写的麻烦。

\texttt{delta\_entries} 列族承担热点目录的增量写入。它的目标不是长期保存完整状态，而是承接高频的 \texttt{mtime}/\texttt{ctime}/\texttt{nlink} 变化，避免每次目录修改都去重写父目录 inode 基值。因为这一列族里的条目短小、生命周期也短，还会被后台折叠线程定期消费，所以这里刻意关掉了压缩，省掉热路径上的编解码开销。

\texttt{system} 列族只放 inode 分配器这类低频系统元信息。单独隔离出来是为了避免它和热点目录项共用一套压缩和缓存行为。访问频率虽然远低于其他列族，但它仍然要和普通元数据一起走原子提交，所以留在同一个 RocksDB 实例里，没有额外搞一套外部存储。

归纳一下，这套列族划分走的是“按访问模式组织布局”的思路，而不是“按语义模块机械拆分”。对本文的讨论来说，列族设计本身就是元数据关键路径优化的一部分，所以得先把“为什么这么分”讲清楚，再讨论“每类记录怎么编码”。

\subsubsection{键编码与目录映射}

在列族边界确定之后，接下来要解决的问题是：目录树中的 inode 和目录项，如何映射为 RocksDB 中可比较、可扫描、可点查的字节串键。RucksFS 对所有多字节整数统一采用大端序编码，以保证数值顺序与字节序顺序一致；不同类型的键再通过单字节前缀区分，其中 \texttt{[I]} 表示 inode 属性记录，\texttt{[D]} 表示目录项记录，\texttt{[X]} 表示增量记录，\texttt{[L]} 表示数据位置记录，\texttt{[S]} 表示符号链接目标记录。在此基础上，表~\ref{tab:key_encoding} 给出几类核心键的具体编码格式。

\begin{table}[htbp]
    \centering
    \caption{核心键编码格式}
    \label{tab:key_encoding}
    \small
    \setlength{\tabcolsep}{4pt}
    \begin{tabular}{p{2.0cm}p{4.6cm}p{4.2cm}}
    \toprule
    \textbf{键类型} & \textbf{编码格式} & \textbf{说明} \\
    \midrule
    inode 键 & \texttt{[I][inode: u64 BE]} & 存放 inode 基础属性 \\
    目录项键 & \texttt{[D][parent: u64 BE][name: UTF-8]} & 存放父目录到子名字的映射 \\
    增量键 & \texttt{[X][inode: u64 BE][seq: u64 BE]} & 存放时间戳或 \texttt{nlink} 增量 \\
    数据位置键 & \texttt{[L][inode: u64 BE]} & 存放 \texttt{DataLocation} \\
    符号链接键 & \texttt{[S][inode: u64 BE]} & 存放符号链接目标路径 \\
    \bottomrule
    \end{tabular}
\end{table}

图~\ref{fig:key_mapping} 以路径 \texttt{/docs/paper.tex} 为例，展示了目录树如何映射到 RocksDB 键空间。核心思想不是存“整条路径”，而是存“父目录到子名字”的边关系。

\begin{figure}[htbp]
    \centering
    \begin{tikzpicture}[
        font=\small,
        node distance=8mm and 13mm,
        tree/.style={draw, rounded corners=2pt, fill=black!5, align=center, minimum width=2.2cm, minimum height=0.8cm},
        kv/.style={draw, fill=black!2, align=left, minimum width=5.9cm, minimum height=0.8cm},
        arr/.style={->, >=stealth, thick}
    ]
    \node[tree] (root) {root\\inode 1};
    \node[tree, below=10mm of root] (docs) {docs\\inode 17};
    \node[tree, below=10mm of docs] (paper) {paper.tex\\inode 42};
    \draw[arr] (root) -- (docs);
    \draw[arr] (docs) -- (paper);

    \node[kv, right=28mm of root] (k1) {\texttt{[D][1][“docs”] -> [17][DIR]}};
    \node[kv, below=7mm of k1] (k2) {\texttt{[I][17] -> inode attrs of docs}};
    \node[kv, below=7mm of k2] (k3) {\texttt{[D][17][“paper.tex”] -> [42][REG]}};
    \node[kv, below=7mm of k3] (k4) {\texttt{[I][42] -> inode attrs of paper.tex}};

    \draw[arr] ($(root.east)!0.5!(docs.east)$) -- (k1.west);
    \draw[arr] ($(docs.east)!0.5!(paper.east)$) -- (k3.west);
    \end{tikzpicture}
    \caption{路径树到 RocksDB 键空间的映射示意}
    \label{fig:key_mapping}
\end{figure}

这种设计遵循的是”以边为中心”的命名空间建模思路。目录树中的每一条父子边，都对应 \texttt{dir\_entries} 列族中的一条记录。\texttt{lookup(parent, name)} 退化为点查一条目录项键；\texttt{readdir(parent)} 退化为扫描同一前缀下的全部目录项键；\texttt{rename} 则退化为删除旧边、插入新边。与使用全路径作为键相比，这种建模方式避免了跨目录重命名时对子树全量重写。第~\ref{sec:pjdfstest}~节中 pjdfstest 的 \texttt{rename}、\texttt{readdir}、\texttt{lookup} 用例将直接检验这一建模方式的正确性。

\textbf{与 TableFS 键编码方案的对比。} 最早将 KV 存储应用于文件系统元数据的 TableFS（Ren et al., 2013）同样采用“父 inode + 文件名”作为目录项键。但 TableFS 使用原生字节序编码 inode 编号，导致同一父目录下的所有目录项在 LSM-tree 的键空间中不一定物理相邻——因为 LevelDB 按字节序比较键值，inode 编号 256 的字节序表示为 \texttt{\textbackslash x01\textbackslash x00\textbackslash x00\textbackslash x01}，而 inode 编号 257 表示为 \texttt{\textbackslash x01\textbackslash x00\textbackslash x00\textbackslash x02}，两者之间可能隔着大量其他 inode 的键。RucksFS 采用大端序编码，使得 inode 编号的大小顺序与字典序一致，同一父目录下的所有子项天然在键空间中连续排列。这一差异在 \texttt{readdir} 场景下尤为关键：TableFS 需要在扫描时额外过滤父目录编号，而 RucksFS 只需一次前缀迭代即可获得全部目录项。此外，大端序也使得前缀 Bloom 过滤器在同一父目录内的假阳性率更低。

\subsubsection{Value 编码与定长序列化}

在实现层，inode 的基础属性由 \texttt{FileAttr} 抽象表示。它对应上层接口最常使用的一组对象属性，也是 \texttt{InodeValue} 编码时需要覆盖的核心字段，其内容主要包括 inode 编号、文件大小、模式位、链接计数以及访问、修改和状态变更时间戳。

在键编码之外，不同类型 value 的布局也经过了针对性设计。inode value 采用 57 字节的定长二进制布局，依次保存版本号、inode 编号、文件大小、模式位、链接计数以及三类时间戳；目录项 value 固定为 12 字节，即 \texttt{[child\_inode: u64 BE][child\_mode: u32 BE]}；delta value 则采用极短的 \texttt{DeltaOp} 编码。

这三种 value 的设计分别服务于不同目标。首先，把 \texttt{mode} 一并放入目录项 value，可以让 \texttt{readdir} 在不额外回查 inode 的情况下直接得到对象类型。其次，delta value 使用短小的 \texttt{DeltaOp} 编码，是为了让热点目录上的属性更新尽可能退化为小对象追加写。最后，inode value 采用定长布局，则有利于在点查路径上减少解析开销。

MetadataServer 在 RocksDB 中存储的 inode 基础值并不使用通用序列化框架，而是采用固定长度的 \texttt{InodeValue} 手工编码。当前布局总长度为 57 字节，包含版本号、inode 编号、文件大小、模式位、链接计数以及三类时间戳。采用固定长度编码有两个直接好处：其一，读写路径只需要按偏移切片，不必承担变长解码开销；其二，调试时可以直接从 RocksDB 中观察二进制布局，问题定位更直接。

固定长度编码也意味着模式演进需要显式处理。为此，\texttt{InodeValue} 的第一个字节保留为格式版本号。只要未来确实需要扩展字段，就可以通过版本字段引入兼容逻辑，而不必打破既有键空间。

\subsection{事务模型与元数据一致性}
\label{sec:txn_model}

\subsubsection{事务边界与原子提交}

RucksFS 的大多数元数据操作都不是单键更新。拿 \texttt{create} 来说，它至少要同时改目录项、目标 inode、数据位置以及父目录属性；\texttt{rename} 往往同时涉及源目录、目标目录、被移动对象和潜在被覆盖对象。所以事务模型要先回答的不是“怎么加锁”，而是“哪些修改必须作为一个不可分割的整体提交”。

本文的做法是不把 RocksDB 原生事务接口直接散在业务代码里调，而是通过 \texttt{StorageBundle} 和 \texttt{AtomicWriteBatch} 把跨列族写入封装起来。调用方先用 \texttt{StorageBundle::begin\_write()} 开一个事务上下文，往里面追加 \texttt{BatchOp}，最后统一提交。这样做的好处是，上层逻辑可以围绕“目录项在不在”、“应该写哪些 inode 记录”、“父目录要追加哪些增量”来组织动作，不用显式处理不同列族的底层细节。提交时底层最终会调到 RocksDB 事务的 \texttt{commit()}，把跨 \texttt{inodes}、\texttt{dir\_entries}、\texttt{delta\_entries} 的修改作为一个原子单元持久化。

\subsubsection{悲观并发控制与冲突处理}

光有原子提交还不够，并发目录操作的正确性也要管。这里用的是 RocksDB TransactionDB 的悲观并发控制（PCC）接口，关键操作是 \texttt{get\_for\_update}：事务在读取目录项键或 inode 键的同时拿排他锁，“检查条件”和“执行修改”就落在同一个事务视图里。对文件系统来说这件事比较重要——\texttt{create}、\texttt{rmdir}、\texttt{rename} 都有典型的 TOCTOU 风险，如果先在事务外检查、再进事务里改，就会留下竞态窗口。

多对象操作还会带来死锁。本文的做法是：凡是可能涉及多个 inode 的事务，都先把相关 inode 编号收集起来排序，再按升序逐个加锁。\texttt{rename} 是最典型的例子——源目录、目标目录、被移动对象、潜在被覆盖对象都可能同时进同一个事务。统一的加锁顺序把“先锁谁、后锁谁”变成一条确定规则，不同线程就不会因为加锁顺序不一致而形成循环等待。当前实现还给锁等待设了 5 秒超时；如果事务冲突一时消不掉，MetadataServer 会先在服务端有限重试，最多 3 次，并做指数退避；重试后还不行，再作为 \texttt{TransactionConflict} 往上返回，在 FUSE 层映射成 \texttt{EAGAIN}。

\subsubsection{增量更新与读写一致性}

事务并不能解决所有热点问题。目录修改里还有一个更隐蔽的瓶颈：父目录属性更新。按 POSIX 语义，\texttt{create}、\texttt{unlink}、\texttt{mkdir}、\texttt{rmdir} 都会动到父目录的 \texttt{mtime}/\texttt{ctime}，涉及子目录的操作还会动 \texttt{nlink}。如果每次都直接读出父目录 inode 再写回，热点目录下冲突就会集中到同一条基值记录上，LSM-tree 擅长的顺序追加反而被“读-改-写”盖过去了。

本文的做法是把这类更新从“基值重写”改成“增量追加”。写路径不直接改父目录在 \texttt{inodes} 里的基础值，而是把变化编码成一条 \texttt{DeltaOp}，追加到 \texttt{delta\_entries} 里。目前的增量类型有 \texttt{IncrementNlink}、\texttt{SetMtime}、\texttt{SetCtime}、\texttt{SetAtime}。普通文件的创建和删除只影响父目录时间戳；\texttt{mkdir}、\texttt{rmdir} 和跨目录移动目录时还会额外引入 \texttt{nlink} 增量。这样热点目录上的高频更新就从“反复覆盖同一条 inode 基值”变成了“往独立的增量记录里追加”。

读取路径这一侧则要负责把“基础值 + 增量”重新合成成最新视图。MetadataServer 在加载 inode 时，优先从 \texttt{InodeFoldedCache} 读已经折叠好的结果；缓存未命中就先读 \texttt{inodes} 里的基础值，再扫属于这个 inode 的 \texttt{DeltaOp}，按顺序叠起来。为了不让读路径的折叠成本无限增长，后台的 \texttt{DeltaCompactionWorker} 会周期性检查 dirty inode，增量条目超过阈值时把折叠后的新基值写回 \texttt{inodes}，同时把已消费的增量删掉。当前默认阈值 32 条、检查周期 5 秒。思路可以概括为：用一点受控的读路径折叠成本，换写路径上的冲突削减。和 Mantle 的 Delta Record 相比，本文的 DeltaOp 沿用了“以增量换写放大”的基本思路，但把追加、折叠和提交都限制在同一个 RocksDB 实例里完成，实现上更轻一些。

这套机制的收益在低并发下看不出来：父目录上的写争抢必须到一定密度，覆盖式更新和增量追加的差距才会拉开。第~\ref{sec:delta-vs-nodelta}~节会用 Delta 和 NoDelta 两种编译变体做对照实验来量化这一点——共享父目录的元数据负载下，关掉 DeltaOp 的构建在客户端数从 32 增长到 64 时出现 40\% 吞吐下滑，而开启 DeltaOp 的构建保持近线性。

\subsubsection{删除语义与句柄生命周期}

事务原子性之外，元数据一致性还体现在删除语义上。POSIX 允许目录项已经从命名空间里消失，已打开的句柄仍然能继续访问原文件。只靠一次 \texttt{unlink} 事务搞不定这件事，因为“名字空间可见性”和“数据回收时机”并不总在同一时刻。本文因此在 MetadataServer 里维护了 \texttt{open\_handles} 和 \texttt{pending\_deletes} 两类状态。\texttt{open()} 成功后对应 inode 的打开计数加一；\texttt{release()} 时减一。

当 \texttt{unlink} 或覆盖式 \texttt{rename} 把某个普通文件的 \texttt{nlink} 降到 0 时，MetadataServer 不会一律立刻删数据。如果这个 inode 还有打开句柄，系统只把它加进 \texttt{pending\_deletes}；等最后一个句柄关掉的时候，再通知客户端到 DataServer 侧真正清理数据。这样目录项删除、链接计数归零、物理数据回收就分成了三个阶段：名字先从命名空间里消失，可达性由 \texttt{nlink} 表示，物理删除则拖到没有进程还持有这个对象时再做。外面看起来，这就对应了 POSIX 要求的“名字已删除、但已有句柄仍然有效”的语义。pjdfstest 里的 \texttt{unlink::open\_file\_not\_freed} 用例直接检验了这一行为，结果见第~\ref{sec:pjdfstest-semantic}~节。

\subsection{核心操作实现}
\label{sec:posix_impl}

RucksFS 的大多数元数据操作走同一条主线：先定位少量相关的目录项键和 inode 记录，在同一事务里完成检查和修改，按需给父目录追加属性增量。这一节挨个看这些操作的流程，也顺便把前面键编码、PCC、DeltaOp 三件事是怎么组合复用的串一下。

\subsubsection{create 与 mkdir}

\texttt{create} 的难点不是“写入一个新 inode”，而是要同时满足三件事：目标名字原本不存在，新对象以完整状态出现，父目录的属性更新和前两件保持原子一致。要是把这些拆成多个独立步骤，就很容易出现中间态——比如 inode 已经建好但目录项还没插进去，或者同名检查通过后被别的并发事务抢先插入。算法~\ref{alg:create} 给出当前实现的 \texttt{create} 流程。

\begin{algorithm}[htbp]
    \caption{\texttt{create(parent, name, mode, uid, gid)} 的核心流程}
    \label{alg:create}
    \begin{algorithmic}[1]
        \State 开启事务 batch
        \State 对目录项键 \texttt{D[parent, name]} 执行 \texttt{get\_for\_update}
        \If{目录项已存在}
            \State 返回 \textsc{EEXIST}
        \EndIf
        \State 分配新 inode 编号，构造普通文件的 \texttt{InodeValue}
        \State 向 \texttt{inodes} 列族写入新 inode 基值
        \State 向 \texttt{dir\_entries} 列族写入父目录到新文件的映射
        \State 向 \texttt{inodes} 列族写入该 inode 的 \texttt{DataLocation}
        \State 向 \texttt{delta\_entries} 追加父目录的 \texttt{SetMtime} 与 \texttt{SetCtime}
        \State 提交事务
        \State 更新本地缓存并返回新文件属性
    \end{algorithmic}
\end{algorithm}

从键空间组织的角度看，\texttt{create} 里能看到边中心建模的直接好处：不用改整条路径，也不用搬子树，只要在 \texttt{dir\_entries} 里插一条新边，再在 \texttt{inodes} 里把目标对象的基础记录和数据位置记录补齐。新对象的可见性是由一条目录项记录加一组 inode 记录共同建立的，不是由哪个单一结构承担。

另一方面，父目录属性不走“覆盖写回基础 inode”那条路，而是追加两条 \texttt{DeltaOp}——一条 \texttt{SetMtime}、一条 \texttt{SetCtime}。这一点把 \texttt{create} 从“目录项插入 + 父目录基值重写”的组合，换成了“目录项插入 + 若干新记录写入 + 增量追加”的组合。前者在热点目录下容易把竞争集中到父目录 inode 上，后者更贴合 LSM-tree 擅长的追加写模式。

\texttt{mkdir} 和 \texttt{create} 共用同一条事务骨架，但比普通文件创建多了一层目录相关的语义约束。新对象的类型位是目录，初始 \texttt{nlink} 为 2；父目录除了更新时间戳之外，还要追加 \texttt{IncrementNlink(+1)}，因为多了一个子目录。普通文件创建不会改父目录 \texttt{nlink}、目录创建会——这点看着细，但直接决定了后面 \texttt{rename}、\texttt{rmdir} 和一致性检查的正确性。

\subsubsection{unlink、rmdir 与 release}

如果说 \texttt{create} 是“建一条新边”，那 \texttt{unlink} 和 \texttt{rmdir} 就分别对应“删一条普通文件边”和“删一条目录边”。共同点是两者都要在事务里同时处理目录项可见性、目标对象状态和父目录属性变化；区别来自 POSIX 对普通文件和目录有不同约束。

\texttt{unlink} 的核心工作流是：锁定目标目录项，确认它不是目录，删掉这条目录项记录，把目标 inode 的 \texttt{nlink} 减一，再按链接计数和打开句柄的状态决定要不要走延迟删除分支。这里的关键点是目录项删除和对象状态更新在同一个事务里，不会出现“名字已经消失但链接计数还没更新”这种中间态。父目录的 \texttt{mtime}/\texttt{ctime} 同样不直接重写基值，继续走增量日志。

\texttt{rmdir} 的难点在空目录检查。如果写成“先扫一遍目录、确认为空、再开事务删除”，检查结果在进事务前就可能失效。本文的做法是把“是不是空”的判断也放进事务里，用一致性视图下的前缀扫描完成。如果目录非空，直接返回 \texttt{ENOTEMPTY}；如果为空，就删掉目标目录项和目标 inode，给父目录追加 \texttt{IncrementNlink(-1)} 加一组时间戳增量。空目录约束是事务本身保证的，不依赖上层约定。

\texttt{release} 补上了删除语义的最后一环。\texttt{unlink} 和覆盖式 \texttt{rename} 可以让名字从命名空间里消失，但只要还有进程持有打开句柄，物理数据就不能马上清。每次客户端关文件句柄时 \texttt{release} 都会把打开计数减一；要是某个 inode 之前已经进了 \texttt{pending\_deletes}、此时计数又刚好降到 0，MetadataServer 才会把这个 inode 返回给客户端，通知它到 DataServer 侧做真正的数据删除。这样“名字可见性”、“元数据可达性”、“物理数据回收”就清晰地分成三个阶段。

\subsubsection{rename}

\texttt{rename} 是前面几处设计选择交汇最密的操作。它既依赖“以边为中心”的键建模，也依赖统一的加锁顺序、覆盖删除语义和父目录属性增量，可以看成整条元数据路径里最能体现系统设计取舍的一个操作。它可能同时牵涉源目录、目标目录、被移动对象、潜在被覆盖对象这四类状态，并且要求整个过程对外原子可见。算法~\ref{alg:rename} 概括了这一流程。

\begin{algorithm}[htbp]
    \caption{\texttt{rename(old\_parent, old\_name, new\_parent, new\_name)} 的核心流程}
    \label{alg:rename}
    \begin{algorithmic}[1]
        \State 读取并锁定源目录项
        \State 若目标目录项存在，则读取其 inode 信息
        \State 收集全部相关 inode 编号并排序，按序执行 \texttt{get\_for\_update}
        \State 检查源对象与目标对象的类型兼容性
        \If{目标存在且为非空目录}
            \State 返回 \textsc{ENOTEMPTY}
        \EndIf
        \If{目标存在}
            \State 按覆盖删除逻辑处理目标对象
        \EndIf
        \State 删除源目录项，写入目标目录项
        \State 更新被移动对象的 \texttt{ctime}
        \If{移动的是目录且跨父目录}
            \State 对旧父目录追加 \texttt{IncrementNlink(-1)}
            \State 对新父目录追加 \texttt{IncrementNlink(+1)}
        \EndIf
        \State 对涉及的父目录追加时间戳增量
        \State 提交事务并返回被覆盖对象列表
    \end{algorithmic}
\end{algorithm}

从实现角度看，\texttt{rename} 的主体可以理解为“删旧边、插新边，按需处理被覆盖对象”。这是边中心建模的直接好处：哪怕是跨目录移动，系统只需要动少量和源路径、目标路径直接相关的记录，不用像全路径键那样对子树做级联改写。\texttt{rename} 的复杂性主要来自语义约束，而不是来自底层数据搬移。

排序加锁是 \texttt{rename} 能避免死锁的关键，不过它的意义不止“避免互相等”。更重要的是，统一的加锁顺序让源目录、目标目录、被覆盖对象能被纳入同一个稳定事务里，这样“检查目标是否存在”、“检查目标目录是否为空”、“删旧目录项”、“写新目录项”、“处理覆盖删除”这一串动作都能保持同一个一致性视图。在外部观察者看来，\texttt{rename} 就是一次原子切换，不会暴露“旧名字已删、新名字未建”这种中间态。

还有个容易忽略的细节：只有目录跨父目录移动时才需要调整父目录 \texttt{nlink}。移动普通文件不会影响任何目录的 \texttt{nlink}；同一父目录内对目录重命名也不会影响父目录 \texttt{nlink}。这些差异直接决定了增量日志里要追加哪些 \texttt{DeltaOp}，也说明 \texttt{rename} 的正确性不仅仅取决于“目录项有没有改对”，还取决于相关目录属性有没有被同步维护。

\subsubsection{readdir、open 与写入协调}

\texttt{readdir} 最能体现键空间局部性。给定父目录 inode 之后，MetadataServer 构造前缀 \texttt{[D][parent: u64 BE]}，用 RocksDB 的前缀迭代器就能顺序扫出这个目录下的所有 dentry。因为目录项 value 里已经带着子对象的 \texttt{mode}，普通 \texttt{readdir} 不用再为每个条目回查一次 inode 就能拿到文件类型。这也回头说明前面在目录项 value 里多存一个 \texttt{child\_mode} 不是冗余，而是专门用来降低目录遍历放大成本的。

\texttt{open} 是元数据面和数据面的衔接点。它本身不做数据读写，而是负责把“这个 inode 的数据该去哪台 DataServer 读写”这一信息交给客户端。所以 \texttt{open} 返回的不只是句柄，还包括对应的 \texttt{DataLocation}。这样后续的 \texttt{read}、\texttt{write}、\texttt{truncate} 就不用每次都回 MetadataServer 重新问一次路由，可以直接访问正确的数据后端。

写入完成之后的元数据更新也体现了 RucksFS 的职责划分：DataServer 只管把字节写到正确位置，并不主动维护文件长度或 \texttt{mtime}；这些语义统一由 MetadataServer 掌握。客户端在 \texttt{write\_data} 成功之后会显式调 \texttt{report\_write(inode, new\_size, mtime)} 去回写元数据。好处是边界清楚：数据路径保持简单，元数据语义集中维护；代价就是客户端要多做一次协调动作。更完整的路由和服务编排细节在下一节展开。
\subsection{数据服务与客户端协调}
\label{sec:data_path}

到上一节为止，讨论的重点还停留在“单个元数据操作在 MetadataServer 内部怎么走”。但真正完整的系统里，用户触发的是一条跨 FUSE 客户端、MetadataServer、DataServer 的路径。文件系统的正确性不光取决于元数据事务本身，还取决于客户端怎么在几个服务之间编排请求、怎么维护句柄到数据位置的映射、怎么把 FUSE 的可见语义落到分离的元数据面和数据面上。这一节因此换成系统级视角，看看这条跨服务路径是怎么搭起来的。

\subsubsection{VfsCore 路由与语义编排}

客户端内部的协调层是 \texttt{VfsCore}。它一方面对 FUSE 暴露统一的 \texttt{VfsOps} 接口，另一方面同时持有 \texttt{MetadataOps} 和 \texttt{DataOps}，负责决定一次请求应该落到哪个服务、要不要跨服务串联、返回结果后还得补哪些后续动作。\texttt{VfsCore} 不是一个简单的 RPC 转发器，而是整条路径的编排者。

请求大致分三类。一类是纯元数据操作——\texttt{lookup}、\texttt{create}、\texttt{mkdir}、\texttt{unlink}、\texttt{rmdir}、\texttt{rename}、\texttt{readdir}、\texttt{getattr}、\texttt{setattr}——直接交给 MetadataServer。一类是纯数据操作——\texttt{read}、\texttt{write}、\texttt{truncate}、\texttt{fsync}——文件已经打开的情况下可以按句柄里缓存的 \texttt{server\_id} 直接路由到对应的 DataServer。剩下一类是跨服务复合操作，比如 \texttt{open} 要先从 MetadataServer 拿到 \texttt{DataLocation}，\texttt{release}、覆盖式 \texttt{rename}、可能触发截断的 \texttt{setattr} 则要在元数据结果返回后再由客户端补一个数据侧的清理或截断请求。

这样安排背后的意图很明确：只涉及名字空间和 inode 状态的语义都约束在 MetadataServer 里，真正写字节、截数据、清数据的工作都留在 DataServer 里，跨两侧的复合语义交给客户端显式协调。结果就是边界比较清楚：MetadataServer 管元数据事务，DataServer 管字节读写，\texttt{VfsCore} 把两边缝成用户看到的 POSIX 行为。代价是客户端会比传统本地文件系统里那种薄代理复杂一些，这是分离式架构必须付出的一层成本。

\subsubsection{DataServer 与原始磁盘后端}

数据侧本文走了一条比较简单的路径。DataServer 只是 \texttt{DataStore} 的服务化封装层，当前唯一的后端是 \texttt{RawDiskDataStore}。它把所有 inode 的内容映射到同一个 \texttt{data.raw} 文件，用固定公式
\[
\text{disk\_offset} = \text{inode} \times \text{max\_file\_size} + \text{file\_offset}
\]
算物理偏移，再用 \texttt{pread}/\texttt{pwrite} 读写。DataServer 不维护复杂的块映射结构，而是把 inode 编号直接当作数据文件里的一个逻辑槽位。

这种固定地址映射的好处挺直接：数据路径不需要空闲空间位图、extent 树或块分配器，寻址是常数时间，服务端状态也很简单。对本文聚焦的“元数据关键路径优化”来说还有个附加好处——数据侧逻辑越简单，第~\ref{sec:experiments}~章里的性能差异就越容易归因到元数据组织、客户端协调和网络通信上，不会被复杂的块管理逻辑搅浑。

但这也是本文要诚实交代的一个短板。每个 inode 都预留 \texttt{max\_file\_size} 大小的逻辑地址空间，目前默认 128\,MiB，大量小文件下空间利用率就会很低。\texttt{delete\_data} 当前还是空操作，对象删除之后物理空间不会立即回收。之所以留着这套方案，不是因为空间管理不重要，而是本文在工程上有意识地把这些复杂度推到后续扩展中，换当前原型在数据路径上的可解释性和实现简洁性。

并发这一块，写入偏移由调用方自己指定，不同客户端往同一 inode 的不同偏移写，就不会因为 DataServer 内部再分配地址而产生额外竞争；\texttt{pread}/\texttt{pwrite} 也保证每次 I/O 以调用为边界完成。如果不同写入请求本身覆盖了重叠区域，那属于上层应用的并发语义问题，不是 DataServer 要额外裁决的状态冲突。

\subsubsection{FUSE 可见语义与缓存策略}

用户态文件系统里，很多“性能表现”首先体现为“内核缓存怎么看这个文件系统”。所以在进实验之前，先得把 FUSE 层的几个可见语义讲清楚，不然第~\ref{sec:experiments}~章里的数字很容易被误读。本文主要控制三件事：目录项缓存、属性缓存、目录遍历接口的选择。

目录项缓存（entry cache）决定内核相信“名字到 inode”的映射在多长时间内有效。本文在 \texttt{lookup} 响应里设 1 秒的 entry TTL，免得目录变更被内核缓存挡太久。属性缓存（attr cache）决定 \texttt{stat}、\texttt{getattr} 等请求的放大程度，这里默认同样用 1 秒 TTL，在一致性和 RPC 次数之间折中。目录遍历接口上显式启用 \texttt{readdirplus}，让目录遍历在返回目录项的同时带上基本属性，省掉“先 \texttt{readdir} 再对每个条目逐个 \texttt{getattr}”的往返。

除了缓存以外，FUSE 语义里还有两个和分离式架构比较相关的点。一是 \texttt{open} 在本文里不光产生句柄，还负责把数据路由信息交给客户端。二是 \texttt{fsync} 目前只保证数据侧 \texttt{flush}，元数据仍依赖 RocksDB 的事务提交和 WAL 持久化，没有做严格的双阶段刷写。也就是说，本文对 FUSE 的实现不是把内核接口原封不动搬到远端服务，而是根据元数据面/数据面分离后的职责边界，给每个接口一个能落地的语义解释。

\subsubsection{gRPC 接口与模块边界}

在这条路径上，gRPC 不只是通信细节，它也在决定“哪些状态由哪个服务直接返回”、“哪些后续动作由客户端继续协调”。所以接口设计本身也是系统结构的一部分。本文在接口层面遵循两条原则：元数据操作和数据操作严格分开；一次 RPC 尽量返回下一阶段需要的全部信息，避免把简单操作拆成多轮往返。

据此，MetadataServer 的接口大体分成四组：命名空间操作（\texttt{Lookup}、\texttt{Create}、\texttt{Rename}、\texttt{Readdir} 等），属性操作（\texttt{GetAttr}、\texttt{SetAttr}），句柄相关（\texttt{Open}、\texttt{Release}），以及写入后的元数据协调（\texttt{ReportWrite}）。DataServer 的接口刻意收敛，只暴露 \texttt{Read}、\texttt{Write}、\texttt{Truncate}、\texttt{DeleteData} 这类直接对字节内容的操作。这样“谁负责对象语义、谁负责数据内容”在接口层面就被固定下来了。

消息格式也服务于同一个目标。元数据响应通常直接带上完整的 \texttt{FileAttr}、\texttt{DataLocation}、\texttt{purged\_inodes} 或 \texttt{needs\_truncate}，省得客户端为了完成一个上层操作还得再发一轮探测 RPC。连接层面，当前原型用的是“每客户端单 gRPC 通道、共享 HTTP/2 连接”，主要为了控制实现复杂度。这肯定不是最终形态，但对本文关注的元数据语义验证和原型级性能评估够用了。

\texttt{VfsCore}、RawDiskDataStore、FUSE 缓存设置、gRPC 接口这几部分合起来组成了 RucksFS 的跨服务执行路径。上一节回答的是“单个操作在元数据层如何正确完成”，这一节补上的是“这些操作在完整系统里如何被协调并暴露给内核和应用”。

\subsection{错误处理与恢复机制}

一旦系统从单进程存储引擎扩展到“客户端 + MetadataServer + DataServer”这种三段式路径，正确性问题就不光来自事务冲突，还来自 RPC 失败、服务崩溃和跨服务状态不一致。因为第~\ref{sec:experiments}~章主要评估正确性和性能结果，不再解释系统行为本身，这里需要先把故障模型和恢复边界讲清楚。

\subsubsection{事务冲突、RPC 失败与错误码映射}

最常见的一类失败是 MetadataServer 内部的事务冲突。对这种情况，本文没选择把 RocksDB 的底层错误直接抛给用户，而是先在服务端有限重试：最多 3 次、指数退避。只有重试还不行才作为 \texttt{TransactionConflict} 往上返回，再在 FUSE 层映射成 \texttt{EAGAIN}。这样偶发冲突就不会直接变成用户可见错误，错误语义也能保持清晰：\texttt{EAGAIN} 表示的是“事务冲突还没消解”，而不是任意 I/O 异常。

另一类失败来自普通的 RPC 或 I/O 调用。这块当前原型选了比较保守的做法：客户端按 gRPC 返回的错误状态向上报相应的文件系统错误，而不会在数据路径上维护待确认写队列，也没有做基于 \texttt{inode + offset} 的幂等重发。说白了，当前只做了“事务冲突自动重试”，并没有做“任意跨服务失败都自动恢复”。这一边界必须说清楚，不然读者会误以为系统已经具备更强的故障屏蔽能力。

\subsubsection{服务崩溃后的恢复语义}

MetadataServer 挂掉的时候，它持有的 RocksDB TransactionDB 会依赖 WAL 做恢复：已经提交但还没刷入 MemTable 的内容会在重启后恢复，没完成提交的事务不会伪装成“已经成功”。也就是说，只要某次元数据操作已经返回成功给客户端，它对应的持久化结果就不会因为 MetadataServer 进程挂掉而丢。客户端侧会把崩溃表现为一次 RPC 失败或重试；必要时，\texttt{open} 这类操作可以重新找 MetadataServer 拿新的 \texttt{DataLocation}。

DataServer 的崩溃语义更简单，也更受当前实现边界约束。如果负责某个文件的 DataServer 不可达，对应的 \texttt{read} 或 \texttt{write} 会直接向上返 \texttt{EIO}。因为当前原型没引入数据副本或多 DataServer 故障切换，数据服务故障不会被透明屏蔽，而是直接暴露为“这个文件暂时访问不了”。又因为 \texttt{delete\_data} 当前仍是空操作，删除路径也就不存在“删除已提交但物理空间没及时回收”这种额外的恢复问题。

\subsubsection{当前实现边界与论文讨论范围}

讲完正常路径和故障路径之后，需要明确一下本文的讨论边界。当前 RucksFS 有意保持“单 MetadataServer + 单默认 DataServer”这条主路径。虽然 \texttt{DataLocation} 和 \texttt{VfsCore::with\_data\_servers()} 已经为多 DataServer 路由预留了接口，但元数据分片、多副本一致性、物理空间回收和严格的双阶段 \texttt{fsync} 都还没做进来。

这不是简单意义上的“没做完”，而是本文研究目标的一部分。本文更关心元数据面和数据面分离之后，如何通过键编码、事务控制、客户端协调和增量更新来优化元数据关键路径，而不是去讨论一个完整的大规模分布式控制面。所以第~\ref{sec:experiments}~章的实验结果也就主要验证这套设计在当前边界内的正确性和性能特征，并不打算证明系统已经覆盖多副本容错、跨节点故障转移或空间回收这种更大范围的工程能力。

\subsection{本章小结}

本章围绕 inode 与 dentry 这两个核心对象展开。元数据按访问模式被拆到不同的列族里，目录树用“以边为中心”的键编码映射到 RocksDB 键空间；事务部分用悲观锁加统一加锁顺序保证并发正确性，父目录属性则通过 DeltaOp 增量记录减少热点冲突。\texttt{create}、\texttt{mkdir}、\texttt{unlink}、\texttt{rmdir}、\texttt{rename}、\texttt{readdir}、\texttt{open} 等具体操作基本都在这套键编码、事务和 DeltaOp 的组合下落地。\texttt{VfsCore} 的路由层、\texttt{RawDiskDataStore} 的简化实现、FUSE 的缓存设置以及 gRPC 接口边界，则把上述单点机制串成一条完整的跨服务执行路径。

需要说明的是，目前的实现只覆盖了这条路径的主干：单 MetadataServer、单默认 DataServer，没有副本，也没有物理空间回收，DataServer 本身是最小可用实现。这是本文在时间和精力上的取舍，把重点放在了“元数据路径在分布式部署下如何组织”这一核心问题上。后续第~\ref{sec:experiments}~章的实验也围绕这条主干展开，通过 pjdfstest 合规性、mdtest 多系统对比和 Delta/NoDelta 消融，检验这些设计是否真的达到了预期。

